\documentclass[12pt, a4paper]{article}
\usepackage[utf8]{inputenc}
\usepackage[brazil]{babel}
\usepackage{booktabs}
\usepackage{longtable}
\usepackage{graphicx}
\usepackage{geometry}
\usepackage{hyperref}
\usepackage{float}
\geometry{margin=2.5cm}

% --- DADOS DO GRUPO ---
\newcommand{\projetotitulo}{Introdução à Manufatura Aditiva}

\begin{document}

\begin{center}
    {\Large\bfseries Termo de Abertura\sprintnum}\\
    {\large \projetotitulo}\\
\end{center}

\section{Identificação do Projeto}

O Projeto a ser desenvolvido é um minicurso de manufatura aditiva (modelamento e impressão 3D) que será ofertado no campus e tem como público alvo estudantes de escolas públicas do ensino médio, mas não limitante a este. \\

\noindent \textbf{Integrantes do Grupo e Papéis Scrum:}
\begin{itemize}
    \item \textbf{Alunos de Ensino Médio} --- Product Owner
    \item \textbf{Yuri Mendonça} --- Developer
    \item \textbf{Lucas Leão} --- Tech Lead
    \item \textbf{Vinícius Vilela Novelo} --- Divulgation / Developer
    \item \textbf{Paulo Arthur} --- Developer
    \item \textbf{Leonardo Oliveira} --- Divulgation / Developer
    \item \textbf{Lucas Moura} --- Scrum Master / Developer
    \item \textbf{Isaque Victor} --- Developer
    \item \textbf{Calouro (A definir)} --- Developer
\end{itemize}

\section{Justificativa}
Este projeto de extensão surge com o propósito de democratizar o acesso à manufatura aditiva, popularmente conhecida como impressão 3D, levando esse conhecimento técnico de forma acessível para a comunidade externa à universidade. Muitas vezes, essa tecnologia é vista pelo público em geral como algo distante, complexo e restrito apenas a indústrias de ponta ou laboratórios de engenharia. Nossa proposta é quebrar essa barreira e demonstrar que a impressão 3D é uma ferramenta perfeitamente viável e extremamente útil para o cotidiano. O foco é capacitar pessoas da comunidade nos princípios básicos de modelagem, utilizando ferramentas como o Solid Edge, e no processo de fatiamento e impressão, provando que qualquer pessoa pode aprender a materializar suas próprias ideias.

O problema real que buscamos endereçar é a exclusão tecnológica e a falta de ferramentas acessíveis para que a comunidade resolva desafios físicos do seu dia a dia. Comumente, o cidadão se depara com situações onde uma peça de plástico de um eletrodoméstico quebra e não é mais fabricada, ou precisa de um suporte específico para casa, um adaptador customizado, ou até mesmo a criação de uma ferramenta simples para auxiliar no trabalho. Sem o acesso à manufatura aditiva, a solução costuma ser o descarte de objetos que poderiam ser consertados ou a dependência de soluções improvisadas. Ao levarmos a impressão 3D para esse público, oferecemos uma alternativa barata, rápida e feita sob medida.

Portanto, a relevância deste projeto está em empoderar a comunidade através da cultura \textit{maker} (faça você mesmo). O objetivo é fornecer uma nova habilidade prática para o público externo. Ao dar a eles o conhecimento para projetar e fabricar suas próprias peças, promovemos a independência, estimulamos a criatividade na resolução de problemas caseiros e mostramos que a tecnologia universitária pode, de fato, melhorar e facilitar a vida das pessoas de forma direta e palpável.

\section{Objetivo Geral e Objetivos Específicos}
\subsection{Objetivo Geral}
Capacitar os alunos convidados no uso de tecnologias de manufatura aditiva e modelagem 3D, aplicando esse conhecimento técnico na criação de soluções práticas e protótipos funcionais para demandas cotidianas até o final do semestre letivo.

\subsection{Objetivos Específicos}
\begin{itemize}
    \item Específico: Capacitar os estudantes no uso do SolidEdge para modelamento 3D;
    \item Mensurável: Avaliar o correto desenvolvimento dos modelos 3D propostos aos alunos;
    \item Atingível: Garantir que cada aluno consiga desenvolver pelo menos uma peça na impressora 3D;
    \item Relevante: Avaliar o impacto dos conhecimentos aprendidos no cotidiano ou pretensões futuras dos estudantes e da comunidade.
\end{itemize}

\section{Escopo Preliminar}
\subsection{Itens inclusos no curso}
\begin{itemize}
    \item Fornecer conhecimentos básicos em desenho técnico;
    \item Disponibilizar o espaço e ferramentas adequadas (Softwares e Impressoras 3D) durante o tempo do curso;
    \item Capacitar os calouros e alunos no modelamento 3D com SolidEdge;
    \item Capacitar os calouros e alunos na impressão 3D;
    \item Possibilitar o desenvolvimento de protótipos de interesse do estudante durante o tempo do curso;
    \item Fornecer documentações informativas a respeito da modelagem 3D no SolidEdge e impressão 3D de pequenas peças.
\end{itemize}
\subsection{Itens exclusos no curso}
\begin{itemize}
    \item Fornecer transporte para os alunos inscritos no curso;
    \item Disponibilizar o espaço e ferramentas adequadas (Softwares e Impressoras 3D) em horários diferentes e/ou além do tempo do curso;
    \item Capacitar demais pessoas não inscritas no curso no modelamento 3D com SolidEdge;
    \item Capacitar demais pessoas não inscritas no curso na impressão 3D;
    \item Possibilitar o desenvolvimento de protótipos de interesse do estudante em horários diferentes e/ou além do tempo do curso;
\end{itemize}

\section{Stakeholders}

Os \textit{stakeholders} (partes interessadas) identificados para este projeto de extensão, bem como seus respectivos papéis, interesses e níveis de influência sobre os resultados, estão descritos na tabela abaixo.

\begin{table}[h!]
    \centering
    \renewcommand{\arraystretch}{1.3} 
    \begin{tabular}{ll p{6cm} c}
        \toprule
        \textbf{Nome / Grupo} & \textbf{Tipo} & \textbf{Interesse} & \textbf{Nível de Influência} \\
        \midrule
        
        Grupo Manufatura Aditiva & Interno & Executar o projeto com sucesso, aplicar conhecimentos técnicos e obter aprovação na disciplina. & Alto \\
        \addlinespace
        
       Virgil Del Duca Almeida & Interno & Garantir o alinhamento pedagógico, o cumprimento dos prazos e a qualidade do material produzido. & Alto \\
        \addlinespace
        
        Comunidade externa local & Externo & Aprender sobre modelagem 3D e fatiamento para aplicar em soluções do seu dia a dia. & Alto \\
        \addlinespace
        
        IFMG Campus Betim & Externo & Viabilizar o espaço para o curso e receber capacitação técnica para a comunidade local. & Médio \\
        \addlinespace
        
        IFMG & Interno & Promover impacto social positivo na comunidade e cumprir seu papel de extensão acadêmica. & Baixo \\
        
        \bottomrule
    \end{tabular}
    \caption{Mapeamento dos Stakeholders do Projeto}
\end{table}

\section{Cronograma macro}

\begin{table}[H]
\centering
\begin{tabular}{l c p{6cm} p{4cm}}
\toprule
\textbf{Sprint} & \textbf{Semanas} & \textbf{Foco Principal} & \textbf{Marco de Entrega} \\
\midrule

Sprint 0 & 1--2 & Kickoff, diagnóstico, formação dos grupos & Marco 1: Termo de Abertura \\

Sprint 1 & 3--4 & Backlog, capacitação e planejamento & Marco 2: Product Backlog + Plano \\

Sprint 2 & 5--6 & Preparação das aulas, criação de slides e início da divulgação & Relatório de Sprint 2 \\

Sprint 3 & 7--8 & Coleta de dados dos alunos, definição de espaço e equipamentos & Relatório de Sprint 3 \\

Revisão & 9 & Sprint Review + Retrospectiva de meio de semestre & Apresentação interna \\

Sprint 4 & 10--11 & Realização das primeiras 2 aulas & Relatório de Sprint 4 \\

Sprint 5 & 12--13 & Realização das últimas 2 aulas e coleta de dados de satisfação dos alunos & Relatório de Sprint 5 \\

Sprint 6 & 14--15 & Finalização e documentação & Marco 3: Relatório Final \\

Encerramento & 16 & Feira de Extensão (apresentação pública) & Marco 4: Apresentação Final \\

\bottomrule
\end{tabular}
\caption{Planejamento de Sprints}
\end{table}

\section{Riscos iniciais}

\subsection{Risco 1: Falhas de impressão devido às propriedades do filamento ABS}
\begin{itemize}
    \item \textbf{Probabilidade:} Alta
    \item \textbf{Impacto:} Médio
    \item \textbf{Descrição:} O filamento ABS exige um controle térmico rigoroso. Há um risco considerável de as peças sofrerem contração durante o resfriamento, resultando em descolamento da mesa de impressão (warping) ou rachaduras entre as camadas depositadas, gerando perda de material e necessidade de reimpressão.
    \item \textbf{Estratégias de Mitigação:} Incluir no escopo das aulas práticas orientações específicas sobre os parâmetros de fatiamento (como a obrigatoriedade do uso de brim ou raft). Além disso, utilizar impressoras 3D com gabinete fechado para manter a temperatura do ambiente controlada e aplicar adesivos de fixação na mesa de impressão.
\end{itemize}

\subsection{Risco 2: Gargalo logístico e filas de espera na etapa de impressão}
\begin{itemize}
    \item \textbf{Probabilidade:} Média
    \item \textbf{Impacto:} Alto
    \item \textbf{Descrição:} O processo de deposição FDM é naturalmente lento. Se os alunos projetarem peças muito grandes ou volumosas, ou se todos finalizarem a etapa de desenho no SolidEdge ao mesmo tempo, haverá uma fila de impressão que pode ultrapassar o prazo do curso, impedindo que todos vivenciem a etapa final de ``testar'' a peça.
    \item \textbf{Estratégias de Mitigação:} Estabelecer regras rígidas de dimensionamento e volume máximo para as peças já na fase de desenho. Configurar parâmetros obrigatórios de baixo preenchimento interno (infill) para agilizar a fabricação de protótipos e organizar um cronograma logístico de ``fornadas'', agrupando várias peças de alunos diferentes em uma mesma impressão.
\end{itemize}

\subsection{Risco 3: Curva de aprendizado prolongada no software SolidEdge}
\begin{itemize}
    \item \textbf{Probabilidade:} Baixa
    \item \textbf{Impacto:} Baixo
    \item \textbf{Descrição:} Como a turma pode ter diferentes níveis de familiaridade com computação e desenho técnico, alguns alunos podem apresentar lentidão para dominar o SolidEdge, atrasando o início da sua etapa de impressão em relação ao restante da turma.
    \item \textbf{Estratégias de Mitigação:} Disponibilizar tutoriais complementares de apoio (vídeos ou apostilas curtas) focados apenas nas ferramentas básicas que serão efetivamente usadas no projeto. Promover o trabalho em pares ou grupos nas primeiras aulas para que os alunos com maior facilidade auxiliem os colegas.
\end{itemize}

\end{document}